\documentclass[11pt,a4paper]{article}
\usepackage[margin=25mm]{geometry}
\usepackage{amsmath,amssymb,amsthm,hyperref}
\setlength{\emergencystretch}{2em}
\newtheorem{lemma}{Lemma}
\newtheorem{theorem}[lemma]{Theorem}
\newtheorem{corollary}[lemma]{Corollary}
\newcommand{\gen}[1]{\langle#1\rangle}
\title{Strict row exits and critical internal orbit budgets}
\author{Human Compression project: L044 continuation}
\date{}
\begin{document}\maketitle
For the first two results, retain the complete ambient string-action
hypotheses of \emph{Uniform prefix growth in an internal pair equality}.
No intersection property or parity assumption is needed there. The
later critical-budget application additionally uses the accepted
source-head reduction to $n=2r$; that reduction retains its smaller-degree
alternating induction and all supporting proofs.

\begin{lemma}[The additional fixed point works before the last label]
\label{extra}
Retain every hypothesis of \emph{Uniform prefix growth in an internal
pair equality}. In particular $O$ has $2q$ points, $q\ge3$,
$k\notin F$, $K=G_{>k}$ is intransitive on $O$, and
$|F|=q-1$. Let $a=\min F$, $x=\max F=a+q-2$, and let
$C\subseteq\operatorname{Fix}(K)$ have size $h\ge1$.
If the whole $H=G_{\ge k}$ fixes $C$ pointwise, then
\[
       |\operatorname{Fix}(G_{>x})|\ge(a-k+2)q+h+1.       \tag{1}
\]
For the two source heads $(k,h)=(2,2)$ and $(5,7)$ this is at
least $2x+3$.
\end{lemma}
\begin{proof}
Use the accepted prefix closure $Z=O^{G_{\le a-2}}$.
It consists of nontrivial $q$-point orbits of
$Q=\gen{\rho_j:j\in F}$ and has at least $(a-k+2)q$ points.
The later tail fixes $Z$, and $Z$ is disjoint from $C$.

The full set $\operatorname{Fix}(K)$ is preserved by
$P=G_{<k}$, because $P$ centralizes $K$. If it were exactly $C$,
then both $P$ and $H$ would preserve the proper nonempty set $C$,
contradicting transitivity of $G=\gen{P,H}$. There is therefore
another $K$-fixed point outside $C$. It is $Q$-fixed and so lies
outside $Z$ as well. Since $G_{>x}\le K$, it is also fixed by
the later tail. This proves (1) for every $x$, not just the last.

Write $a=k+1+e$ with $e\ge0$. Subtracting $2x+1$ from (1)
gives $e(q-2)+q+h-2k+2$. This is at least three for $(2,2)$
and at least two for $(5,7)$. In particular (1) is at least
$2x+3$ in both cases.
\end{proof}

\begin{corollary}[The saturated nonfirst orbit gives the rank bound]
\label{exclude}
In either source head of Lemma~\ref{extra}, suppose every colour
after $x$ has two ambient-eligible edges. Then
\[
                        r\le\lfloor(n-1)/2\rfloor.
\]
\end{corollary}
\begin{proof}
If $x=r-1$, (1) gives $n\ge2x+3=2r+1$.
Otherwise put $d=r-1-x>0$ and $T=G_{>x}$.
The nonempty head $C$ is fixed by $K$ and hence by $T$.
As $T$ has positive rank, it is intransitive on $\Omega$.
The accepted ambient-interval theorem, applied to this full
terminal interval, gives
\[
           2d\le |\operatorname{Supp}(T)|-c,
\]
where $c\ge1$ counts its nontrivial orbits. Hence its support
has at least $2d+1$ points, even when it is transitive on that
support. Together with Lemma~\ref{extra}, this yields
$n\ge2x+3+2d+1=2r+2$, which is stronger than required.
All selected tail edges use the ambient omissions: a local
paired condition alone would not justify this step.
\end{proof}

\begin{theorem}[The exact critical orbit budget]\label{budget}
For this theorem let $G\le A_n$ have connected string C-marking,
all maximal omissions intransitive, and assume $n=2r$ under the
accepted critical source-head reduction. Fix $k$ and use the
normalized assignments of
\emph{Ambient transport and internal pair deficits}: every tail
colour is assigned exactly once, with two ambient-eligible edges
in its assigned actual $H=G_{\ge k}$-orbit, and every assigned
orbit has intransitive full local $K=G_{>k}$.
Let $f=|\operatorname{Fix}(H)|$, let $c$ count the nontrivial
$H$-orbits, and let $p$ be the orbit assigned colour $k$.
Let an ambient involution $\tau$ centralize $H$ and lie in every
omission $G_j$ for $j\ge k$.

For the singleton head $k=2$, $f\ge2$, suppose also that $H$
is intransitive on its support, as provided by the accepted
critical odd-tail reduction. Suppose the actual two-point head is fixed by $H$.
If $\tau$ acts oddly inside an actual $H$-orbit with nonempty
assigned set, this critical configuration is impossible.

For the six-point head $k=5$, suppose $H$ fixes its specified
seven-point head. Internal odd action on any orbit other than $p$
with nonempty assigned set is impossible. If it acts oddly on $p$,
only the following two exact budgets remain:
\[
\begin{array}{c|c|l}
f&c&\text{orbit sizes and deficits}\\\hline
8&1&|O_p|=2q,\quad d_p=-1,\\
7&2&|O_p|=2q,\quad d_p=-1,\quad
 |O_s|\text{ odd},\quad d_s=-\tfrac12.
\end{array}                                                     \tag{2}
\]
Here $q$ is odd and at least three; the second orbit in the
second row has odd size at least five. On $O_p$, the pair images
of $F_p=\{k,\ldots,k+q-2\}$ form the ordinary path, and every
other tail pair image is trivial. Signed lifts are not classified.
These two rows are reductions, not exclusions.
\end{theorem}
\begin{proof}
Write $d_s=|F_s|-|O_s|/2$ and $S=\sum_s d_s$.
The exact assignment partition and $n=2r$ give
\[
                       S=f/2-k.                         \tag{3}
\]
The normalized deficit theorem gives $d_p\le0$ and
$d_s\le-\tfrac12$ for $s\ne p$. An internal odd central
involution is free on its orbit, which therefore has $2q$ points
with $q$ odd. If $q=1$, no ambient-eligible assigned colour can
lie there, because its two vertices are already joined by
$\tau\in G_j$. Thus a first-assigned internal orbit has $q\ge3$.
The internal pair-forest theorem gives $d_s\le-1$ in every
internal nonempty orbit; for an empty one of size $2q$ its
exact deficit is $-q$.

If $p$ is internal, then
$S\le-1-(c-1)/2$, so (3) gives $f+c\le2k-1$.
For the singleton head $f\ge2$ and $c\ge2$ contradict this.
For the six-point head $f\ge7$ leaves $f+c\le9$.
Since $O_p$ is even and $n$ is even, the sole-orbit case has
$f$ even, forcing $(f,c)=(8,1)$. Otherwise only $(7,2)$ is
possible, and the second orbit is odd. Equality in each bound
forces precisely the deficits in (2). A three-point second orbit
cannot support its required one assigned colour, since an involution
on three points has at most one edge. Its size is therefore at least
five. The accepted internal equality theorem gives consecutive pair
path labels; since $k\in F_p$, they start at $k$.

Now suppose an orbit $s\ne p$ is internal and odd. The assertion assumes its assigned set is nonempty.
For a nonempty assignment $q\ge3$ and $d_s$ is an integer.
If $d_s\le-2$, then $S\le-2$, contradicting (3), since
$f\ge2k-3$ for both heads. Thus $d_s=-1$.
The complete internal equality theorem gives exactly the
nonfirst-orbit hypotheses of Lemma~\ref{extra}; the actual fixed
head supplies its $H$-fixed set $C$ of size two or seven.
Every later colour retains its ambient eligible pair by the
assignment hypothesis. Corollary~\ref{exclude} is the contradiction.
\end{proof}

An empty internal odd orbit of size at least six is also excluded:
its exact deficit is at most $-3$, whereas (3) is at least $-3/2$.
An empty two-point orbit remains a separate source case.

\begin{lemma}[The exceptional head enters a centralizer-fixed point]
\label{head}
Use the actual six-point source head
$D=\{1,\ldots,6\}$, with $\gen{\rho_0,\rho_1,\rho_2}$
transitive and supported on $D$. Put $\tau=\rho_3$,
$\sigma=\rho_4$, and suppose
\[
 \tau=(6\ b)\delta,\qquad
 \operatorname{Supp}(\delta)\subseteq\Omega\setminus(D\cup\{b\}),
 \qquad \sigma|_D=1.
\]
Retain full ambient IP and the accepted natural alternating tail
\[
              G_{\ge3}=A_Z,\qquad
              Z=\Omega\setminus\{1,\ldots,5\}.
\]
If $x=b^\sigma\ne b$, then $x^\tau=x$.
\end{lemma}
\begin{proof}
Let $A=\gen{\rho_0,\rho_1,\rho_2}$.
As $A$ is transitive on $D$, the conjugates of $\tau$ by $A$
include $(v\ b)\delta$ for every $v\in D$. Products of two
such conjugates cancel $\delta$ and give the pure three-cycles
on $\{b,v,w\}$. Hence $G_{\le3}$ contains the natural
$A_{D\cup\{b\}}$, fixing every point outside that set.

The point $x$ lies outside $D\cup\{b\}$, since $\sigma$
fixes $D$ and $x\ne b$. Conjugating this pure alternating group
by $\sigma$ and taking the generated group shows that $G_{\le4}$
contains pure $A_{D\cup\{b,x\}}$. We use the elementary
alternating-overlap fact here: two natural alternating groups
on sets of size at least four with at least two common points
generate the natural alternating group on the union. Their
three-cycles within the two sets, conjugated using the common
points, give the three-cycles crossing the two sets, proving the fact.

If $x^\tau\ne x$, conjugate once more by $\tau\in G_{\le4}$.
The set $D\cup\{b\}$ is $\tau$-invariant, so this gives pure
$A_{D\cup\{b,x,x^\tau\}}\le G_{\le4}$ by the same overlap fact.
In particular it contains pure $A_{\{6,b,x,x^\tau\}}$.
All four points lie in $Z$, so this last group also lies in
$G_{\ge3}=A_Z$. Full IP puts it in
$G_{\le4}\cap G_{\ge3}=\gen{\tau,\sigma}$.
But a group generated by two involutions is dihedral, and all
its odd-order elements belong to its cyclic rotation subgroup.
The two three-cycles on $\{6,b,x\}$ and $\{6,b,x^\tau\}$ do
not commute. This contradiction proves the claim.
\end{proof}

\begin{corollary}[The remaining first-assigned budget has an odd anchor]
In the actual six-point head, additionally use the source interface
that $B=G_{\ge4}=\gen{\sigma,H}$ is transitive on
$\Omega\setminus D$, that $H=G_{\ge5}$ fixes $D\cup\{b\}$,
and that $\tau$ centralizes $H$. Under the first-assigned internal
odd case in Theorem~\ref{budget}, the first row of (2) is impossible.
In its second row,
\[
             b^\sigma\in O_s\cap\operatorname{Fix}(G_{>5}),
\]
where $O_s$ is the odd orbit of size at least five. The involution
$\tau$ fixes that entire orbit pointwise.
\end{corollary}
\begin{proof}
Transitivity of $B$ makes $x=b^\sigma\ne b$.
If $f=8$ and $c=1$, the $H$-fixed points outside $D$ are exactly
$b$ and one further point $z$. If $x=z$, the set $\{b,z\}$ is
preserved by both $\sigma$ and $H$, contradicting transitivity
of $B$. Thus $x\in O_p$, where $\tau$ is free. This contradicts
Lemma~\ref{head}.

If $(f,c)=(7,2)$, no $H$-fixed point outside $D\cup\{b\}$
exists, so $x$ is in one of the two nontrivial orbits.
Lemma~\ref{head} excludes $O_p$. The odd orbit $O_s$ is
$\tau$-invariant: it cannot be exchanged with the even orbit
$O_p$, and commutation maps actual $H$-orbits to actual orbits.
An involution centralizing a transitive action on an odd number
of points fixes that orbit pointwise. Finally $G_{>5}$ fixes
$b$ and commutes with $\sigma=\rho_4$, hence fixes $x$.
\end{proof}

\paragraph{Source and limits.}
The target is the all-paired part of source Proposition 5.31,
physical pages 42--43, in Cameron--Fernandes--Leemans--Mixer,
\emph{Highest rank of a polytope for $A_n$}, arXiv:1605.09173v1.
The orbit-deficit strategy and fixed-support target are borrowed.
The additional fixed point before the last label, ambient terminal
closure and critical-budget routing are the present arguments.
The pure alternating overlap and dihedral obstruction reuse the
mechanism of the accepted dihedral attachment proofs, here with the
actual six-point head and its natural alternating tail.
The accepted odd-tail reduction supplies $n=2r$ only under its
complete standing hypotheses, which remain paid. No ladder shape
for the first-assigned internal orbit is presumed.
The generic budgets (2) have not been excluded without the
additional literal head interface. With that interface, only the
odd-anchor row remains. This residual, the empty two-point case,
assignment existence outside the stated ambient boundary, finer
source diagrams and whole L044 remain outstanding. No computation, formal result,
original credit or question credit is added.
\end{document}
